\documentclass[a4paper, 12pt, oneside, headspline]{scrartcl}
% !TeX root = ../main.tex
% Variablen welche innerhalb der gesamten Arbeit zur Verfügung stehen sollen
\newcommand{\titleDocument}{Bachelorarbeit}
\newcommand{\subjectDocument}{Test-driven development mit Python}
\newcommand{\institution}{Hochschule für angewandte Wissenschaften Augsburg}
\newcommand{\insSection}{Fakultät für Informatik}
\newcommand{\name}{Maximilian Konter}
\newcommand{\email}{maximilian.konter@hs-augsburg.de}

% Farbige links
\usepackage{color}

% Umlaute unter UTF8 nutzen
\usepackage[utf8]{inputenc}

% deutsche Silbentrennung
\usepackage[ngerman]{babel}

% Grafiken aus PNG Dateien einbinden
\usepackage{graphicx}

% bricht lange URLs "schoend" um
\usepackage[hyphens,obeyspaces,spaces]{url}

% mehrseitige Tabellen ermöglichen
\usepackage{longtable}

% für Tabellen
\usepackage{array}

% Schaltet den zusätzlichen Zwischenraum ab, den LaTeX normalerweise nach einem Satzzeichen einfügt.
\frenchspacing

% Paket für Zeilenabstand
\usepackage{setspace}

% für Bildbezeichner
\usepackage{capt-of}

% für Stichwortverzeichnis
\usepackage{makeidx}

% Für Glossar
\usepackage{glossaries}

% Bessere Listen
\usepackage{enumitem}

%\bibliographystyle{}
\usepackage[
    maxnames=6,
    minnames=1,
    style=apa,
    backend=biber,
    bibencoding=utf8,
    sortlocale=de_DE,
]{biblatex}
\DeclareLanguageMapping{ngerman}{ngerman-apa}
\DefineBibliographyStrings{ngerman}{
    andothers = {{et\,al\adddot}},
}

% für Listings
\usepackage{listings}
\lstset{numbers=left, numberstyle=\tiny, numbersep=5pt, keywordstyle=\color{black}\bfseries, stringstyle=\ttfamily,showstringspaces=false,basicstyle=\footnotesize,captionpos=b}
\lstset{language=python}

% Indexerstellung
\makeindex

% Abkürzungsverzeichnis
\usepackage[german]{nomencl}
\let\abbrev\nomenclature

% Abkürzungsverzeichnis LiveTex Version
\renewcommand{\nomname}{Abkürzungsverzeichnis}
\setlength{\nomlabelwidth}{.25\hsize}
\renewcommand{\nomlabel}[1]{#1 \dotfill}
\setlength{\nomitemsep}{-\parsep}
\makenomenclature
\makeglossaries

% Disable single lines at the start of a paragraph (Schusterjungen)
\clubpenalty = 10000
% Disable single lines at the end of a paragraph (Hurenkinder)
\widowpenalty = 10000
\displaywidowpenalty = 10000

% Hyperlinks im Text
\usepackage[bookmarksnumbered,pdftitle={\titleDocument},hyperfootnotes=false]{hyperref}
\hypersetup{colorlinks, citecolor=red, linkcolor=blue, urlcolor=black}

% neue Kopfzeilen mit fancypaket
\usepackage{fancyhdr} %Paket laden
\pagestyle{fancy} %eigener Seitenstil
\fancyhf{} %alle Kopf- und Fußzeilenfelder bereinigen
\fancyhead[L]{\nouppercase{\leftmark}} %Kopfzeile links
\fancyhead[C]{} %zentrierte Kopfzeile
\fancyhead[R]{\thepage} %Kopfzeile rechts
\renewcommand{\headrulewidth}{0.4pt} %obere Trennlinie
\fancyfoot[C]{\thepage} %Seitennummer
\renewcommand{\footrulewidth}{0.4pt} %untere Trennlinie

% Bibliothek hinzufuegen
\addbibresource{bachelor.bib}

\begin{document}
    % Leere Seite am Anfang
    \newpage
    \thispagestyle{empty} % erzeugt Seite ohne Kopf- / Fusszeile
    \section*{ }

    % Titelseite
    % !TeX root = ../bachelor.tex
\thispagestyle{empty}

\begin{figure}[t]
 \centering
 \includegraphics[width=0.3\textwidth]{abb/logoHSA}
\end{figure}


\begin{verbatim}


\end{verbatim}

\begin{center}
\Large{\institution}
\end{center}


\begin{center}
\Large{\insSection}
\end{center}
\begin{verbatim}




\end{verbatim}
\begin{center}
\doublespacing
\textbf{\LARGE{\titleDocument}}
\singlespacing
\begin{verbatim}

\end{verbatim}
\textbf{{~\subjectDocument~}}
\end{center}
\begin{verbatim}

\end{verbatim}
\begin{center}

\end{center}
\begin{verbatim}

\end{verbatim}
\begin{center}
\textbf{zur Erlangung des akademischen Grades \\ Bachelor of Science}
\end{center}
\begin{verbatim}





\end{verbatim}

\begin{footnotesize}
\begin{flushleft}
\begin{tabular}{llll}
\textbf{Thema:} & & \subjectDocument & \\
& & \\
\textbf{Autor:} & & \name & \\
& & \email & \\
& & MatNr. 951004 & \\
& & \\
\textbf{Version vom:} & & \today &\\
& & \\
\textbf{1. Betreuer:} & & Dipl.-Inf. (FH), Dipl.-De Erich Seifert, MA &\\
\textbf{2. BetreuerIn:} & & Prof. Dr.-Ing. Alexandra Teynor &\\
\end{tabular}
\end{flushleft}
\end{footnotesize}

    % 1.5 facher Zeilenabstand
    \onehalfspacing

    % Einleitung / Abstract
    % !TeX root = ../bachelor.tex
\begin{center}
    \vspace{0.9cm}
    \textbf{Abstract}
\end{center}

\begin{abstract}
Da die Anforderungen an Software im Laufe der Zeit stark angestiegen sind,
haben sich einige Technologien entwickelt, die dazu verwendet werden diese
Anforderungen zu gewährleisten. Eine dieser Technologien ist das Test-driven
development. Diese Arbeit befasst sich mit dieser Technologie im Bezug zu den
aktuellen Test-Tools der Programmiersprache Python und analysiert diese auf
ihre Tauglichkeit im Bezug zum Test-driven development.
\newline

Dabei werden zuerst verschiedenen Arten von Tools untersucht, die unit-testing-,
mock-testing- und fuzz-testing Tools. Diese werden anhand ihrer
Anwendbarkeit, Effizient, Komplexität und Erweiterbarkeit beschrieben und
analysiert, um mit Hilfe eines Vergleichs derer einen Leitfaden für das am
besten für Test-driven development geeignete Tool zu bieten. Das daraus
resultierende Ergebnis soll aufzeigen, welche Toolkombinationen die beste Wahl
anhand der vorgegebenen Kriterien darstellen. Diese sind Anwendbarkeit,
Effizienz, Komplexität und Erweiterbarkeit.
\newline

Die dabei entstandenen Ergebnisse zeigen auf, dass nicht alle Tools optimal
für Test-driven development geeignet sind. Schlussendlich ergaben sich vier
Tools, welche gemeinsam eingesetzt, das optimale Toolset für die allgemeine
Anwendung bieten. Dabei ist jedoch zu beachten, dass dies immer vom
Anwendungsfall, sowie den Anforderungen an die Software und an die Tools
abhängig ist. Dabei ergab sich, dass \lstinline{pytest} zusammen mit
\lstinline{mocktest} sowie \lstinline{doctest} und \lstinline{hypothesis}
gemeinsam eine sehr gute Wahl darstellen.
\end{abstract}


    % einfacher Zeilenabstand
    \singlespacing

    % Inhaltsverzeichnis anzeigen
    \newpage
    \tableofcontents

    % das Abbildungsverzeichnis
    \newpage
    % Abbildungsverzeichnis soll im Inhaltsverzeichnis auftauchen
    \addcontentsline{toc}{section}{Abbildungsverzeichnis}
    % Abbildungsverzeichnis endgueltig anzeigen
    \listoffigures

    % das Tabellenverzeichnis
    \newpage
    % Abbildungsverzeichnis soll im Inhaltsverzeichnis auftauchen
    \addcontentsline{toc}{section}{Tabellenverzeichnis}
    % \fancyhead[L]{Abbildungsverzeichnis / Abkürzungsverzeichnis} %Kopfzeile links
    % Abbildungsverzeichnis endgueltig anzeigen
    \listoftables

    % das Listingverzeichnis
    \newpage
    % Listingverzeichnis soll im Inhaltsverzeichnis auftauchen
    \addcontentsline{toc}{section}{Listingverzeichnis}
    \fancyhead[L]{Listingverzeichnis} %Kopfzeile links
    \renewcommand{\lstlistlistingname}{Listingverzeichnis}
    \lstlistoflistings

    % das Abkürzungsverzeichnis
    \newpage
    % Abkürzungsverzeichnis soll im Inhaltsverzeichnis auftauchen
    \addcontentsline{toc}{section}{Abkürzungsverzeichnis}
    % das Abkürzungsverzeichnis entgültige Ausgeben
    \fancyhead[L]{Abkürzungsverzeichnis} %Kopfzeile links
    % !TeX root = ../../main.tex
\nomenclature{GPL}{GNU General Public License}
\nomenclature{GNU}{GNU is not Unix}
\nomenclature{LGPL}{GNU Lesser General Public License}
\nomenclature{TDD}{Test-driven development}
\nomenclature{GUI}{Graphical User Interface}
\nomenclature{STDLIB}{Python Standard Bibliothek}
\nomenclature{VCS}{\gls{vcs}}
    \printnomenclature

    %%%%% Beginn Dokument %%%%%
    \newpage
    \fancyhead[L]{\nouppercase{\leftmark}} %Kopfzeile links

    % 1,5 facher Zeilenabstand
    \onehalfspacing

    % einzelne Kapitel
    % !TeX root = ../bachelor.tex
\section{Einleitung}\label{einleitung}

TDD wird in der heutigen Softwareentwicklung immer verbreiteter und beliebter.
Die Ansprüche an Software sind in den letzten Jahren immer weiter gestiegen.
Dies liegt vor allem an der Reichweite, die Software heute hat. So besitzen
im Jahr 2018 bereits circa 66\% aller Menschen ein Smartphone
(\cite{FraukeSchobelt:Smartphone}). Im Arbeitsleben ist ein PC meist gar nicht
mehr weg zu denken. Doch mit den steigenden Nutzerzahlen steigen auch die
Anforderungen, welche die Nutzer an die Software stellen. Somit wird die Anzahl
der gefundenen \Glspl{bug} dementsprechend größer.

Im weltweiten Markt gibt es viele große Unternehmen, die gegenseitig um die
Nutzer kämpfen. Selbstverständlich präferieren die Nutzer denjenigen Anbieter,
welcher die bessere Software bietet. Dies kann sich heute jedoch stetig ändern.
Mit der steigenden Anzahl an Bugs, die gefunden werden, steigt auch die Anzahl
der Nutzer, die von diesen Bugs betroffen sind. Diese Bugs sollen natürlich
schnellstmöglich gefixt werden um so zu verhindern, dass die Nutzer die Software
wechseln.

So schwer es ist seine Nutzer zu halten, umso schwerer ist es, zum Start einer
Software Nutzer zu akquirieren. Es gibt bereits Software, die ähnliche Services
anbieten. So ist es noch schwerer dem Markt beizutreten. Die Anforderungen
werden durch die bereits am Mark tätigen Firmen gesetzt, damit sollten Fehler,
die bereits gelöst wurden nicht mehr auftauchen.

Für Unternehmen sind diese Anforderungen meist schwer zu meistern, weshalb
Software meist mit Fehlern released wird, um diese dann von den Nutzern
aufdecken zu lassen und zu fixen.

Sowohl als Entwickler als auch als Arbeitgeber muss man sich bei der Wahl der
Programmiersprache Gedanken darüber ob und wie einfach eine Sprache für
Test-driven development ein zu setzen ist. Der wichtigste Aspekt bei diesem
Prozess ist die Auswahl und die Qualität der von der Sprache zur Verfügung
gestellten Tools.
    % !TeX root = ../bachelor.tex
\section{Python Test-Tool Analyse}\label{python-tools}
In diesem Kapitel werden die einzelnen Tools anhand der in Kapitel 
\ref{methodik} beschriebenen Methodik analysiert und verglichen. Am ende des 
Kapitel wird für jede Kategorie (unit-, \gls{mock}- und \gls{fuzz}) eine 
Empfehlung anhand der jeweiligen Analyse gegeben, welches Tool die beste 
Eignung für die Anwendung von TDD hat.

% !TeX root = ../bachelor.tex
\subsection{Tools der Standard Bibliothek}\label{python-tools:stdlib}

Die Standard Bibliothek von Python bietet zwei verschiedene Test-Tools(\cite{wiki.python:PythonTestingToolsTaxonomy}).
Zum einen ist dies
\href{http://pyunit.sourceforge.net/pyunit.html}{unittest}\footnote{http://pyunit.sourceforge.net/pyunit.html}
und zum anderen
\href{https://docs.python.org/3/library/doctest.html}{doctest}\footnote{https://docs.python.org/3/library/doctest.html}.
Diese beiden Tools reichen sind im ihrem Umfang bereits so vielseitig dass, es einfach ist eine hohe Test-Abdeckung eines Programms oder eine Bibliothek zu erreichen.
\newline
\\
Beide Tools zählen zu den "`Unit Testing Tools"'(\cite{wiki.python:PythonTestingToolsTaxonomy}) -
auf Deutsch Modul Test-Tools - mit deren Hilfe die einzelnen Module eines Programms
getestet werden können. In einem Programm oder einer Bibliothek wären dies die einzelnen
Funktionen und Methoden.

% !TeX root = ../bachelor.tex
\subsubsection{unittest}\label{python-tools:unittest}

Unittest

% !TeX root = ../bachelor.tex
\subsubsection{doctest}\label{python-tools:doctest}

"`Das doctest Modul sucht nach Textstücken, die wie interaktive Python-Sitzungen
aussehen, und führt diese Sitzungen dann aus, um sicherzustellen, dass sie genau
wie gezeigt funktionieren."'\footnote{Übersetzt aus dem Englischen}
(\cite{docs.python:doctest}). Diese Textstücke müssen sich in Kommentaren
befinden, da sie sonst von Python als Code interpretiert werden.

\lstinline{doctest} bietet dem Nutzer eine Möglichkeit mithilfe von einem Test
den Code zu dokumentieren. Da \lstinline{doctest} lediglich Code ausführt wie in
einer interaktiven Python-Sitzung, ist auch nur das möglich aus zu führen was im
Code geschrieben ist. Test Fälle oder gar \Glspl{mock} sind hierbei nur bedingt
realisierbar, \Glspl{fixture} hingegen sind teilweise realisierbar in Form von
Code der vor dem Test ausgeführt wird.

\lstinline{doctest} selbst nutzt keine Assert Funktionen oder Methoden,
stattdessen wird der Output eines Befehls überprüft. So ergibt eine Funktion
\lstinline{return_3(s=None, i=None)} bei \lstinline{s=True} eine \lstinline{'3'}
und bei \lstinline{i=True} eine \lstinline{3}. Dieses Beispiel ist in Listing
\ref{listing:doctest:example} zu sehen, um einen Output zu erzeugen wurde noch
ein fehlerhafter Test hinzu gefügt. Dieser Output ist in Listing
\ref{listing:doctest:example_output} zu sehen.

Möchte der Entwickler einen etwas größeren Test schreiben, so kann er sich einer
Funktionalität von \lstinline{doctest} bedienen, die es Ihm ermöglicht Tests in
eine externe Datei zu verlagern. Dabei wird die Datei als ein \Gls{docstring}
behandelt, wodurch sie keine \lstinline{"""} benötigt. Um jedoch Code in diese
Datei ausführen zu können muss das zu testende Modul importiert werden.

Durch das schreiben der Tests in den \Glspl{docstring} werden die Module in
denen Tests geschrieben wurden schnell sehr unübersichtlich und lang, wenn der
Entwickler große Tests schreibt. Zwar wird durch externe Textdateien Abhilfe
geschaffen, dennoch sind zu lange und komplexe doctests schwer zu lesen und
nach zu vollziehen.
\newline

Das Listing \ref{listing:doctest:advanced} zeigt den Code aus
\ref{listing:base:my_module} mit \Glspl{docstring} versehen. Dieser Test wurde
mithilfe der \lstinline{main} Funktion von Python ausgeführt. Der in Listing
\ref{listing:doctest:advanced_text} ausgelagerte Test, wurde mithilfe des CLIs
von \lstinline{doctest} ausgeführt. Da beide Test keine Fehler werfen existiert
auch kein Output für diese Tests.
Die einzige Möglichkeit hier einen Output zu bekommen wäre mit \lstinline{-v} im
CLI oder \lstinline{verbose=True} im Funktionsaufruf. Die dort dargestellten
Informationen sind lediglich welche Kommandos ausgeführt wurden, was erwartet
wurde und das, dass erwartete eingetroffen ist. Ein kleines Beispiel ist in
Listing \ref{listing:doctest:output} zu sehen.

Da \lstinline{doctest} selbst keine Test Fälle unterstützt besitzt
\nameref{python-tools:unittest} eine Integration für \lstinline{doctest},
diese ermöglicht es Tests aus Kommentaren sowie Textdateien in unittest zu
integrieren und zu gliedern. Wie \nameref{python-tools:unittest}, ist auch
doctest in der STDLIB, wodurch keine externen Abhängigkeiten geladen werden
müssen.
\newline

Im Bezug auf die Anwendbarkeit von \lstinline{unittest} lässt sich sagen, dass
nicht alles zur Verfügung steht um TDD zu betreiben. Tests können nicht vor der
eigentlichen Funktionalität geschrieben werden und  \Glspl{fixture} und
\Glspl{mock} werden nicht unterstützt.

Die Anforderungen für \lstinline{doctest} sind sehr gering. Das testen geschieht
mithilfe der interaktiven Python-Sitzung, welche jedem Python Entwickler bekannt
ist. Die Tests selbst sind aufrufe der geschriebenen Funktionen und Methoden und
ein Abgleich des Outputs mit dem Erwarteten Wert. Demnach kann ein Entwickler
sehr Effizient mit \lstinline{doctest} umgehen ohne viel wissen zu benötigen.

\lstinline{doctest} selbst bietet keinerlei Komplexität von sich aus,
stattdessen ist die Komplexität die möglich ist vom Entwickler abhängig, da
dieser die Tests komplett selber schreiben muss. Je nach zu testender
Funktion/Methode kann dies abhängig des Test Aufwands schnell unübersichtlich
werden, wenn viel Code abseits des eigentlichen Tests benötigt wird.

Als Erweiterung bestehen lediglich die Integration in
\lstinline{unittest}\footcite{docs.python:doctest} sowie
\lstinline{pytest}\footcite{docs.pytest.org:4.4}.
\newline
\newline
\\
\textbf{Muss nach Zusammenfassung}
Da sich mit Doctest leider Funktionen nicht präparieren lassen, ist dieses
Test Modul für die alleinige Anwendung in TDD nicht nutzbar. Das Tool bietet
keine Möglichkeiten Objekte zu \gls{mock}en wodurch Tests an nicht
implementieren Methoden und Funktionen scheitern. Auch \Glspl{fixture} sind nur
bedingt realisierbar und benötigen viel Code der dupliziert werden muss, da
dieser vor und nach jedem Test geschrieben werden muss.

% !TeX root = ../bachelor.tex
\subsection{Tools abseits der Standard Bibliothek}\label{python-tools:extlib}
Abseits der Standard Bibliothek gibt es einige Tools deren Nutzung von
Vorteil gegenüber der STDLIB ist. Diese werden unter diesem Punkt aufgeführt.
\newline
\\
Auf \url{https://wiki.python.org/moin/PythonTestingToolsTaxonomy} werden viele
externe Tools gelistet, jedoch scheinen viele inaktiv zu sein da ihre \glspl{commit}
teilweise mehr als ein Jahr zurück liegen. Dies ist weder für eine Bibliothek noch
für ein Tool ein gutes Zeichen, da sich die Anforderungen stetig ändern und niemals
alle Bugs gefixt sind.
\newline
\\
Manche der dort aufgelisteten Tools sind bereits oder werden gerade in andere Tools
integriert. Dies kann zum einem sein, da ein Tool eine Erweiterung für ein anderes
war und die Entwickler die Änderungen angenommen haben und zum anderen um die Tools
zu verbessern und mehr Entwickler zur Verfügung zu haben. Eventuell sind auch
andere Gründe dafür verantwortlich, jedoch war dies der Grund bei zum Beispiel
\href{https://github.com/Yelp/Testify/}{Testify}
von Yelp \footnote{https://github.com/Yelp/Testify/}.
\newline
\\
Da sowohl Software als auch Programmiersprachen sich stetig weiter entwickeln werden in
dieser Arbeit nur jene Tools behandelt die sich diesen Entwicklungen Anpassen, sei dies
durch das unterstützen der aktuellsten Version von Python als auch durch neue Innovationen,
sowie \Gls{bug}-fixes. Aus diesem Grund werden Tools deren letzter \Gls{commit} weiter
als ein Jahr zurück liegt hier nicht behandelt.
\newline
\\
Lässt man zusätzlich die Erweiterungen von Tool zunächst außen vor, so bleiben nach
\cite{wiki.python:PythonTestingToolsTaxonomy} folgende Modul Test-Tools zur Verfügung:
\begin{itemize}
    \item \href{https://github.com/pytest-dev/pytest/}{pytest}\footnote{https://github.com/pytest-dev/pytest/}
\end{itemize}
Tools wie \href{https://www.reahl.org/docs/4.0/devtools/tofu.d.html}{reahl.tofu}\footlabel{reahl.tofu}{https://www.reahl.org/docs/4.0/devtools/tofu.d.html} oder
\href{https://pypi.org/project/zope.testing/}{zope.testing}\footlabel{zope.testing}{https://pypi.org/project/zope.testing/} sind zwar mehr oder weniger aktiv, da sie
beide auf die neusten Python Versionen patchen, jedoch bieten sie sonst keinen Mehrwert in ihren Updates.
\href{https://www.reahl.org/docs/4.0/devtools/tofu.d.html}{reahl.tofu}\footref{reahl.tofu} selbst ist
auch nur eine Erweiterung für bestehende Test Tools, wie zum Beispiel pytest, weshalb es hier nicht behandelt wird. So wird
\href{https://pypi.org/project/zope.testing/}{zope.testing}\footref{zope.testing} auch nicht behandelt, da dieses Tool
wie bereits beschrieben nicht aktiv weiter entwickelt wird und es eher dafür gemacht wurde \href{http://www.zope.org/en/latest/}{zope}\footnote{http://www.zope.org/en/latest/}
Applikationen zu testen und nicht Python Applikationen.

Ein sehr bekanntest Projekt ist \href{https://pypi.org/project/nose/1.3.7/}{nose}\footnote{https://pypi.org/project/nose/1.3.7/}, welches aber
nicht mehr in Entwicklung ist. Der Nachfolger \href{https://pypi.org/project/nose2/}{nose2}\footnote{https://pypi.org/project/nose2/} allerdings
ist unter sehr aktiver Entwicklung und auch wenn das Tool nicht von \cite{wiki.python:PythonTestingToolsTaxonomy} erwähnt wird, so wird es doch
hier mit verglichen werden.

Auch sehr bekannt ist \href{https://pypi.org/project/testtools/}{testtools}\footnote{https://pypi.org/project/testtools/}, jedoch ist auch
hier keine wirkliche weitere Entwicklung zu verzeichnen.
\newline
\\
Demnach wäre die liste an Modul Test-Tools folgende:
\begin{itemize}
    \item \href{https://github.com/pytest-dev/pytest/}{pytest}\footnote{https://github.com/pytest-dev/pytest/}
    \item \href{https://pypi.org/project/nose2/}{nose2}\footnote{https://pypi.org/project/nose2/}
\end{itemize}

Die Quelle aus dem Offiziellen Python Wiki beschreibt sonst keine weiteren Tools.
Sucht auf Google nach weiteren Tools, so findet man die oben gelisteten wenn man
nach den gleichen Kriterien Tools heraus filtert.
\newline
\\
Als weitere Kategorie werden \Gls{mock}-Tools geführt. Auch wenn fast alle
unittest-Tools integriertes \gls{mock}ing haben, so lässt sich mit diesen
Tools meist mehr erreichen. Es wurden die gleichen Filter-Kriterien verwendet wie bei den
Test-Tools oben.
\begin{itemize}
    \item \href{https://www.reahl.org/docs/4.0/devtools/stubble.d.html}{stubble}\footnote{https://www.reahl.org/docs/4.0/devtools/stubble.d.html}
    \item \href{https://github.com/timbertson/mocktest/tree/master}{mocktest}\footnote{https://github.com/timbertson/mocktest/tree/master}
    \item \href{https://github.com/bkabrda/flexmock}{flexmock}\footnote{https://github.com/bkabrda/flexmock}
    \item \href{https://bitbucket.org/DavidVilla/python-doublex}{python-doublex}\footnote{https://bitbucket.org/DavidVilla/python-doublex}
    \item \href{https://github.com/ionelmc/python-aspectlib}{python-aspectlib}\footnote{https://github.com/ionelmc/python-aspectlib}
\end{itemize}
Auch hier finden sich sonst keine weiteren Tools die Open-source sind.
\newline
\\
Als letzte Kategorie werden hier die \Gls{fuzz} Tools behandelt, da diese eine
gute Möglichkeit bieten Code ausgiebig zu testen. Das wohl umfangreichste und nach den obigen Kriterien
einzige Tool ist \href{https://github.com/HypothesisWorks/hypothesis}{hypothesis}\footnote{https://github.com/HypothesisWorks/hypothesis}.

% !TeX root = ../bachelor.tex
\paragraph{pytest}\label{python-tools:pytest}\mbox{}
\newline
Das am 04. August 2009 in Version 1.0.0 veröffentlichte Tool pytest (auch py.test genannt)
ist ein sehr umfangreiches und weit entwickeltes Tool. Seit 2009 wird das Tool stets
weiter entwickelt und vorangetrieben, wodurch es eine Menge an Features gewonnen hat.
\noindent
Die Basis Features von pytest sind folgende:
\begin{itemize}
    \item Simple \lstinline{assert} statements.
    \subitem Kein \lstinline{self.assert}
    \subitem Error Überprüfung mit \Gls{context}
    \item Informativer Output in Farbe
    \subitem Der gesamte Output kann angepasst werden
    \item Feature reiche \Glspl{fixture}
    \subitem Vordefinierte \Glspl{fixture} von pytest
    \subitem Geteilte \Glspl{fixture} unter Tests
    \subitem Globale \Glspl{fixture} zwischen Modulen
    \subitem Parametrisierung von Test als \Glspl{fixture}
    \item Überprüfung von \lstinline{stdout} und \lstinline{stderr}
    \item Gliedern von Test in Fällen
    \subitem ... durch Markierung
    \subitem ... durch Nodes (Auswahl der Modul Abhängigkeit)
    \subitem ... durch String Abgleich der Funktions-Namen
\end{itemize}
\noindent
Da pytest nicht in der STDLIB ist, muss es mit einem \Gls{paket} installiert werden.
Dabei werden für \lstinline{pytest 4.4.0} die in Listing \ref{listing:pytest:requirements}
gezeigten Abhängigkeiten installiert. Demnach benötigt pytest sechs externe Abhängigkeiten
zusätzlich zu sich selbst.
\newline
\\
Wie anhand der Basis Features erkennbar ist, bietet pytest einiges um Tests zu schreiben und
aus zu führen. So ist mit den gebotenen \Glspl{fixture} bereits eine Voraussetzung mehr als
erfüllt, da pytest viele verschiedene Arten bietet \Glspl{fixture} zu benutzen.
\Glspl{fixture} werden in pytest allerdings nicht mit \lstinline{setUp} und \lstinline{tearDown}
geschrieben, sondern werden mithilfe eines \glspl{decorator} markiert und den Test Funktionen
als Parameter übergeben. Um die \lstinline{setUp} und \lstinline{tearDown} funktionalität zu
bekommen muss lediglich das keyword \lstinline{yield} verwendet werden, die \Gls{fixture} wird
dann den Code bis zum \lstinline{yield} ausführen und nach der Funktion den Rest nach
\lstinline{yield} ausführen.
\newline
\\
Diese lassen sich durch zusätzliche Parameter weiter anpassen. Diese
Funktionen sind in Kapitel fünf nach \cite{docs.pytest.org:4.4} beschrieben.

Des weiteren bietet pytest auch \gls{mock}ing, so lässt sich beispielsweise mit
\lstinline{monkeypatch.setattr()} der return Wert einer Funktion ersetzen. Dies wird
in pytest automatisch am ende der Funktion rückgängig gemacht, wodurch der Entwickler
sich mehr auf das eigentliche Testen konzentrieren kann.
\newline
\\
Auch hier bietet pytest weitere Möglichkeiten \Glspl{mock} zu verwenden, dazu hat
\cite{docs.pytest.org:4.4} in Kapitel sieben einige Worte geschrieben.

Ein weiteres sehr praktisches Feature von pytest ist die Gliederung in Fälle. Dies ist bei pytest
sehr fein einstellbar, so lässt sich wie in den Basis Features bereits beschrieben ein Test mit
einer oder mehreren Markierungen versehen wodurch eine Gliederung nach Markierung entsteht. Auch
durch das selektieren bestimmten Wörtern lassen sich Tests nach ihrem Namen gliedern, so entsteht
einerseits eine Gliederung zur Ausführung der Tests und andererseits eine Gliederung für die
Entwickler die sie selbst im Code sehen können. Als letzte alternative lassen sich Tests anhand
ihrer Module ausführen, dies geschieht durch die Angabe der Module. So würde
\lstinline{test_datei.py::TestKlasse::test_methode} den Test \lstinline{test_methode} der Klasse
\lstinline{TestKlasse} in der Datei \lstinline{test_datei.py} ausführen.
\newline
\\
Diese Features werden in Kapitel sechs von \cite{docs.pytest.org:4.4} beschrieben.

Selbstverständlich bietet pytest noch weitere Features jedoch sind diese nicht zwangsläufig notwendig
um TDD zu betreiben und sind mehr ein nice to have Feature als wirklich benötigt. Für eine
Vollständige Auflistung und Erklärung aller Features kann jederzeit unter
\url{https://docs.pytest.org/en/latest/} die aktuellste Version der Dokumentation abgefragt werden.

Demnach bietet pytest alles um TDD anwenden zu können und die sieben zusätzlichen Abhängigkeiten
die ein Entwickler bei der Nutzung von pytest eingeht sollten keinen Entwickler davon abhalten
pytest und seine Features zu genießen.
\newline
\\
Da pytest, wie bereits erwähnt keine neuen assert Methoden hinzufügt lässt sich sehr schnell und
einfach ein Test schreiben. Selbst die Nutzung von \Glspl{fixture} ist in pytest sehr einfach, da
lediglich eine Funktion geschrieben werden muss die mit \lstinline{@pytest.fixture} markiert wurde
und der Test Funktion oder Methode als Parameter übergeben wird. Den Rest erledigt pytest selbst
im Hintergrund. Genauso einfach gestaltet sich die Nutzung von \Glspl{mock}.

Die Effizienz, die bei der Nutzung von pytest entsteht ist demnach sehr hoch. Denn der Entwickler
muss nicht wirklich etwas neues dazu lernen um Tests zu schreiben oder zu präparieren. Genauso
leicht kann ein Entwickler den Output von pytest auswerten, da dieser erstens, in Farbe ist, was
die Lesbarkeit deutlich erhöht, zweitens sehr gut gegliedert ist und Wichtige Strukturen klar
darstellt und drittens nach den Wünschen der Entwickler sich gestalten und verbessern lässt.
\newline
\\
Durch die eben genannten Features im Bezug auf das schreiben von Tests mit assert, sowie
\Glspl{fixture} und \Glspl{mock} lässt sich sagen das pytest sehr viel Funktionalität bietet
wobei es trotzdem Struktur im Code bietet und komplexe Features einfach ermöglicht. Demnach
kann ein Entwickler mit wenig Code viel Testing Funktionalität erstellen.
\newline
\\
Sollten dem Entwickler die Features von pytest nicht ausreichen, so findet man unter
\url{http://plugincompat.herokuapp.com/} eine Liste von 618 (Stand: 2. April 2019) Erweiterungen
für \lstinline{pytest 4.3.0}. pytest selbst kann allerdings auch als Erweiterung zu unittest
genutzt werden, indem ein Entwickler pytest zum ausführen der unittests verwendet. Dadurch bietet
sich dem Entwickler ein verbesserter Output des Test Ergebnisses.
\newline
\\
Im Listing \ref{listing:pytest:advanced} und \ref{listing:pytest:advanced_output} befindet sich der
Code und der Output zum Test von dem in Listing \ref{listing:base:my_module} definierten Modul.
% !TeX root = ../bachelor.tex
\subsubsection{\Gls{mock}ing Tools}\label{python-tools:extlib:mock}

Um eine bessere Gliederung zur Verfügung zu stellen, werden die \Gls{mock}ing
Tools in diesem Kapitel gegliedert.
\\
Wie in \nameref{python-tools} bereits beschrieben, werden hier Tools analysiert,
die eine Erweiterung für Test Tools im Bezug auf verbessertes \gls{mock}ing zur 
Verfügung stellen.

\input{bachelor/3_2_2_1_stubble}
\input{bachelor/3_2_2_2_mocktest}
\input{bachelor/3_2_2_3_flexmock}
\input{bachelor/3_2_2_4_python_doublex}
% !TeX root = ../bachelor.tex
\subsubsection{Fuzz-testing Tools}\label{python-tools:extlib:fuzz}

Zusätzlich zu den hier behandelten Test runnern und den \gls{mock}ing Tools werden
zusätzlich noch Fuzz-testing Tools behandelt. Unter Fuzz-testing versteht man das
Testen eines Programms oder einem Modul mit zufälligen Werten. Diese Werte können
komplett zufällig sein, aber auch einem gewissen Format entsprechen oder einem Typ.

Fuzz-testing funktioniert dabei anders als die gewohnten Methoden des Testens. Der
Normale Flow ist, dass Test Daten präpariert werden, dann werden Tests mit diesen
Daten aus geführt und danach validiert, dass das Ergebnis richtig ist. Mit
Fuzz-testing sieht das ganze etwas anders aus, zuerst wird festgelegt, welchem
Schema die Daten entsprechen müssen (Beispiel: Es wird eine IP erwartet,
\lstinline|[1-9]{1,3}.[1-9]{1,3}.[1-9]{1,3}.[1-9]{1,3}|), danach wird mit den Daten
der Test ausgeführt und danach validiert (\cite{hypothesis:doc:4.18.0}).

Der unterschied ist nun, das der Entwickler selbst sich nicht Daten ausdenken muss.
Im genannten Beispiel müsste der Entwickler sich verschiedene IP Adressen ausdenken,
die entweder richtig oder Falsch sind und danach diese überprüfen. Fuzz-testing
übernimmt das ausdenken der Daten und sorgt dafür, dass diese der Spezifikation
entsprechen. So kann eine viel größere Anzahl an Tests mit erheblich weniger aufwand
auf einem Modul oder Programm ausgeführt werden.

Das ist aber noch nicht alles, schlägt ein Test fehl, mit Daten die per Spezifikation
richtig sind, so wird der Wert der Daten für später gespeichert und sofern möglich
wird versucht mit ähnlichen Daten dieser Fehler zu erreichen. Dadurch hat ein
Entwickler am Ende der Tests Daten die fehlgeschlagen sind. Mit diesen kann er dann
die Weiteren Tests befüllen um so zu verhindern, dass diese Daten zu Fehlern führen.
Im Normalfall übernimmt dies allerdings das Tool, wodurch der Entwickler sich nur
darum kümmern muss, dass die Tests erfolgreich verlaufen.

Von Zeit zu Zeit entsteht dadurch ein Katalog an Daten für verschiedene Tests, der
vorgibt bei welchen Daten ein Test fehl geschlagen ist. Dadurch wird verhindert,
dass beim ändern des Codes dieser Fehler wieder auftritt. Führt man dies über den
gesamten Entwicklungsprozess hinweg durch, so erhält man am Ende ein sehr ausgiebig
getestetes Programm.

Diese Methode des Testens ist auch "`property based testing"', also Eigenschafts-
basierte Tests genannt. In dieser Arbeit, wird allerdings weiterhin der Englische
Begriff Fuzz-testing genutzt.

\input{bachelor/3_2_3_hypothesis}

% !TeX root = ../bachelor.tex
\subsection{Vergleich der Tools}\label{vergleich}

Die in Kapitel \ref{python-tools:stdlib} und \ref{python-tools:extlib}
analysierten Tools, werden anhand der Ergebnisse aus der Anwendung des
jeweiligen Tools in diesem Kapitel verglichen. Der Vergleich wird in unit- und
\gls{mock} Tools gegliedert. Die \gls{fuzz}-testing Tools werden hier
nicht verglichen, da nur ein Tool analysiert wurde.
Am Ende dieses Kapitels wird für die jeweilige Kategorie eine Empfehlung 
ausgesprochen, die anhand der Ergebnisse des Vergleichs festgelegt wird. Anhand 
dieser Empfehlungen werden in Kapitel \ref{kombinierung} verschiedene 
Möglichkeiten aufgezeigt, die Tools untereinander zu kombinieren um mehr 
Funktionalität für eine mögliche Anwendung zu bekommen.

\subsubsection{Unit-testing Tools}\label{vergleich:unit}
Die unit-testing Tools belaufen sich auf drei Tools, \lstinline{unittest},
\lstinline{doctest} und \lstinline{pytest}. Die ersten beiden Tools sind beide
in der STDLIB während \lstinline{pytest} mit sechs weiteren Paketen betrieben
werden muss. Während \lstinline{unittest} und \lstinline{pytest} sich sehr
ähnlich sind, ist \lstinline{doctest} vorkommen anders im Bezug auf Anwendung
und gegebener Funktionalität. \lstinline{doctest} unterstützt keine Testfälle
oder Testumgebungen, stattdessen wird im \gls{docstring} der jeweiligen
Funktion oder Methode der Test geschrieben. Zwar ist es möglich diesen in eine
externe Datei aus zu lagern um so einen Zentralen Punkt zu schaffen an dem sich
alle Tests befinden (Testumgebung/Testfälle), jedoch ist hier die Datei das
Konstrukt, dass die Gliederung ermöglicht. Sowohl \lstinline{unittest} als auch
\lstinline{pytest} bieten hier die Möglichkeiten Tests in Testumgebungen
(Klassen) und Testfällen (Funktionen/Methoden) zu gliedern. \lstinline{pytest}
bietet zusätzlich auch noch die Möglichkeit diese Tests zu markieren um eine
noch feinere Gliederung zu schaffen oder um zwischen Testumgebungen, Testfälle
zu kombinieren. Im Fall der Gliederung und Verwaltung von Tests ist hier
\lstinline{pytest} sowohl \lstinline{unittest} als auch \lstinline{doctest}
voraus. Lässt man die \lstinline{unittest} Integration für \lstinline{doctest}
außen vor, so ist \lstinline{unittest} in diesem Fall der Gewinner. Wird
allerdings die Anwendung der \lstinline{unittest} Integration von
\lstinline{doctest} in Betracht gezogen, so sind \lstinline{unittest} und
\lstinline{doctest} auf einem Level, was die Gliederung und Verwaltung von
Tests betrifft.

Ein Feature, dass eine Gliederung voraussetzt sind die Fixtures. Da
\lstinline{doctest} keine Gliederung bietet sind Fixtures für dieses Tool
nicht als Feature integriert. Es ist zwar möglich Fixtures manuell zu
schreiben, jedoch ist der aufwand, der dafür betrieben werden muss zu hoch, als
dass es sich rentieren würde. Die \lstinline{unittest} Integration ermöglicht es
zwar Fixtures zu nutzen, jedoch wird diese Integration hier wieder außen vor
gelassen. \lstinline{unittest} und \lstinline{pytest} hingegen unterstützen
Testfixtures mit ihren \lstinline{setUp()} und \lstinline{tearDown()} Methoden.
Beide Tools bieten die Möglichkeit die Fixtures sowohl für jedem Testfall, jede
Testumgebung oder jedes Modul aus zu führen. \lstinline{pytest} bietet
zusätzlich noch weitere Funktionalität, sowie fertige Fixtures die dem
Entwickler zur Nutzung zur Verfügung stehen, welche in Kapitel
\ref{python-tools:pytest} beschrieben wurden und in der
\href{https://docs.pytest.org/en/latest/fixture.html}{Dokumentation}\footnote{\url{https://docs.pytest.org/en/latest/fixture.html}}
zu finden sind. Im Bezug auf Testfixtures ist also \lstinline{pytest} das Tool
mit den meisten Features, gefolgt von \lstinline{unittest}. \lstinline{doctest}
scheidet bei dieser Kategorie vollkommen aus, da die Funktionalität gänzlich
fehlt.

Als nächstes integriertes Feature der unit-testing Tools sind \Glspl{mock} und
\Glspl{stub} zu betrachten. Um für TDD in Frage zu kommen sollte ein
unit-testing Tool \Glspl{stub} und \Glspl{mock} unterstützen, eine ausführliche
Erklärung dazu ist in Kapitel \ref{python-tools} zu finden. Wie in Kapitel
\ref{python-tools:doctest} aus der Analyse von \lstinline{doctest} hervorgeht,
gibt es weder die Möglichkeit einen \Gls{mock} noch einen \Gls{stub} auf zu
setzen, weshalb \lstinline{doctest} in dieser Kategorie als unbrauchbar gesehen
werden kann. \lstinline{unittest} als Modul selbst, besitzt keine \gls{mock}ing
Funktionalität, jedoch ist es möglich, dass Submodul \lstinline{unittest.mock}
dafür zu verwenden. \lstinline{unittest.mock} ist zwar eine Erweiterung, jedoch
in Form eines Submoduls, weshalb es zu \lstinline{unittest} dazu gehört und in
diesem Vergleich beachtet wird. Damit unterstützt \lstinline{unittest} sowohl
das \gls{mock}en als auch das \gls{stub}en. \lstinline{pytest} hingegen verfügt
in seiner Basisausführung lediglich über die Möglichkeit des \gls{stub}ens,
nicht jedoch über die Möglichkeit ein Objekt zu \gls{mock}en. In diesem Falls
ist also \lstinline{unittest} der klare Gewinner, gefolgt von
\lstinline{pytest}.
\newline

Im Bezug auf die Effizienz der unit-testing Tools lässt sich folgendes
feststellen. \lstinline{doctest} ist zwar für jeden Entwickler einfach und
verständlich anwendbar, jedoch ist die Vorarbeit um Tests auf zu setzen im
Vergleich zu \lstinline{unittest} und \lstinline{pytest} höher. Sowohl
\lstinline{unittest} als auch \lstinline{pytest} benötigen das gleiche Maß an
Vorarbeit um Tests auf zu setzen, wobei der Aufwand das jeweilige Tool ein zu
setzen bei \lstinline{unittest} etwas höher ist, da es so viele verschiedene
\lstinline{assert} Methoden zu merken gibt. Die Auswertung der Tests erfolgt
bei allen drei unit-testing Tools auf dem Terminal. \lstinline{doctest} und
\lstinline{unittest} haben beide einen Output der voneinander fast nicht zu
trennen ist. Beide Tools haben den Nachteil gegenüber \lstinline{pytest}, dass
ihr Output nicht farbig sonder Monochrom ist.  Aus diesem Grund ist 
\lstinline{pytest}, \lstinline{unittest} und \lstinline{doctest} im Bezug auf
Effiziente Anwendung vor zu ziehen, wobei \lstinline{unittest} noch etwas 
besser als \lstinline{doctest} ist.
\newline

Der Aspekt der Komplexität ist in zwei unter Kategorien gegliedert, einerseits
ist da die Fülle an Funktionen die geboten werden und Komplexität liefern und
zum anderen ist da die Menge an Code die zu schreiben ist und die Anwendung und
Wartung Komplexer gestaltet. Im Falle der ersten Kategorie, erweist sich
\lstinline{doctest} als nicht all zu komplex. Fast die Gesamte gebotene
Funktionalität bezieht sich auf die Ausführung, jedoch nicht das schreiben der
Tests, welches alleine durch die Fähigkeit Python zu schreiben bestimmt wird.
\lstinline{unittest} und \lstinline{pytest} hingegen bieten dem Entwickler
einiges an Funktionalität um Tests zu schreiben und aus zu führen. Auch hier
kann \lstinline{unittest} nur mithalten, wenn das Submodul
\lstinline{unittest.mock} mit in die Wertung aufgenommen wird, da dieses
weitere Funktionalität bietet. \lstinline{unittest} und \lstinline{pytest}
haben viele gemeinsame Features, wovon lediglich das \gls{mock}ing in
\lstinline{pytest} nicht ohne Erweiterungen verfügbar ist. Ansonsten bietet
\lstinline{pytest} viel zusätzliche Funktionalität um Tests feiner ein zu
stellen oder um irgendwelche edge cases ab zu decken. In dieser Kategorie ist
\lstinline{pytest} also etwas besser als \lstinline{unittest} und beide Tools
sind \lstinline{doctest} bei weitem überlegen im Bezug auf die gebotenen
Funktionalitäten.

In der zweiten Kategorie, der Komplexität im Bezug auf Code Lesbarkeit und
Umfang erweisen sich \lstinline{unittest} und \lstinline{pytest} als gleichauf.
Beide Tools haben einen ähnlichen Aufbau, mit den der Code geschrieben wird
(vgl. Listing \ref{listing:unittest:advanced} und
\ref{listing:pytest:advanced}). Die Lesbarkeit ist bei beiden Tools gleich gut
und die Menge an Code die benötigt wird gleicht sich auch. Hier ist es dem
Entwickler überlassen zu entscheiden, welches Tool den schöneren Code schreibt.
\lstinline{doctest} hingegen, produziert viel Code, der mit zunehmender
Funktionalität an Lesbarkeit verliert.Im Bezug auf die Komplexität der Tools
ist \lstinline{pytest} also knapp vor \lstinline{unittest}, gefolgt von
\lstinline{doctest}.
\newline

Als letztes ist die Erweiterbarkeit der unit-testing Tools zu vergleichen.
Beginnend mit \lstinline{doctest}, lässt sich sagen, das es zwei sehr gute
Erweiterungen gibt, die das Tool erheblich verbessern. Diese Erweiterungen sind
die Integration in \lstinline{unittest} und \lstinline{pytest}, erstere wird
von der Python Sprache selbst geboten und ist in \lstinline{unittest} verfügbar
und zweite wird von \lstinline{pytest} geboten und ist auch im Basis Paket
enthalten. Abgesehen von diesen zwei gibt es noch weitere Integrationen in
andere Toolsets, jedoch wurden diese hier aus verschiedenen Gründen nicht
behandelt und werden deshalb nicht erwähnt. Das Modul \lstinline{unittest}
verfügt über einige Erweiterungen welche neue Funktionalität hinzu fügen oder
alte Funktionalität verbessern. Das Tool mit den meisten Erweiterungen ist
\lstinline{pytest}, welches derzeit 618 (Stand: 2. April
2019\footnote{\url{http://plugincompat.herokuapp.com/}}) Erweiterungen
bietet, wovon die meisten sehr nützliche Funktionalität bieten.
\newline

Im Vergleich der unit-testing Tools ergibt sich also ein klarer Gewinner,
\lstinline{pytest}. Das Tool ist in fast allen Analysierten Aspekten die Beste
Wahl für den Einsatz mit TDD. Jedoch bietet \lstinline{unittest} für
Entwickler, die keine Externen Abhängigkeiten möchten eine alternative, die
durchaus mit \lstinline{pytest} mithalten kann.

\subsubsection{Mock-testing Tools}\label{vergleich:mock}
Die analysierten \gls{mock}-testing Tools, sind \lstinline{stubble},
\lstinline{mocktest}, \lstinline{flexmock} und \lstinline{doublex}. Alle vier
bieten verschiedene Möglichkeiten \Glspl{mock} und/oder \Glspl{stub} ein zu
setzen. Da es sich hier um Tools handelt, deren Funktionalität sich um das
\gls{mock}en von Objekten dreht, wird der erste Analysierte Aspekt, die
Anwendbarkeit, auf die gebotene Auswahl an Funktionalität im Bezug auf
\gls{mock}en und \gls{stub}en betrachtet.

Das einzige Tool, das bei der Anwendung auf den zu Testenden Code keine
Probleme mit sich barg, ist \lstinline{mocktest}. Das Tool bietet im Aspekt der
Anwendbarkeit sämtliche Features um effektiv \Glspl{mock} und \Glspl{stub} ein
zu setzen. \lstinline{mocktest} bietet seinem Anwender eine Deskriptive schreib
weise, mit der es problemlos möglich ist Tests zu schreiben. Des weiteren wird
eine Testisolation geboten, mit deren Hilfe es möglich ist Test besser
voneinander ab zu schotten.Zusätzlich benötigt das Tool keine weiteren
Abhängigkeiten zu sich selbst.
Direkt danach wäre \lstinline{stubble} das  nächst beste Tool im Aspekt der
Anwendbarkeit. So bietet das Tool zwar alles um  es für TDD zu verwenden,
jedoch ist das \gls{stub}en einer nicht existierenden Methode nur unter
Einschränkungen Möglich. Die Hauptfeatures von \lstinline{stubble} sind dabei
die vorgefertigten \Glspl{stub}, welche nützliche Features für den Entwickler
bieten. Um einen \Gls{stub} zu erstellen muss die zu \gls{stub}ende Klasse mit
\lstinline{stubble} neu geschrieben werden.Zusätzlich benötigt
\lstinline{stubble} eine weitere Abhängigkeit zu sich selbst.
An dritter Stelle ist das Tool \lstinline{flexmock} welches alles bietet,
abgesehen vom ersetzen einer nicht geschriebenen Methode, um TDD an zu wenden
(vgl. Listing \ref{listing:flexmock:exception}). Wäre dieses Manko nicht, würde
das Tool im Aspekt der Anwendbarkeit sehr weit oben stehen. Flexmock vereint
die \Glspl{mock} und \Glspl{stub} in einer Klasse wodurch die gesamte
Funktionalität mit einem Objekt geboten wird. Die schreib weise der Tests ist
hierbei, wie bei \lstinline{mocktest}, deskriptiv. Auch die benötigten
Abhängigkeiten von \lstinline{flexmock} sprechen für das Tool, da keine
weiteren zu sich selbst benötigt werden. Als letztes Tool im Aspekt der
Anwendbarkeit ist \lstinline{doublex} zu betrachten. Das Tool bietet zwar
einiges an Funktionalität um zu \gls{mock}en oder zu \gls{stub}en, jedoch ist
die Anwendung dieser Features nicht nutzbar für TDD. Es werden vier
verschiedene Interfaces geboten, deren Funktionalität sich überschneidet und
die jeweils ein Objekt hervorbringen, dass ein Duplikat vom Ziel Objekt ist.
Die daraus resultierenden Anwendungszwecke bieten wenig Nutzbarkeit für TDD.
Hinzu kommt, dass \lstinline{doublex} drei weitere Abhängigkeiten benötigt und
anscheinend nicht mehr aktiv weiter entwickelt wird.

Demnach ist im Aspekt der Anwendbarkeit \lstinline{mocktest} also
\lstinline{stubble} und \lstinline{flexmock} vor zu ziehen. \lstinline{doublex}
fällt aufgrund dessen, dass es nicht für TDD geeignet ist und derzeit nicht
aktiv weiter entwickelt wird komplett raus. Aus diesem Grund wird dieses Tool
im weiteren verlauf des Vergleichs nicht mehr berücksichtigt.
\newline

Als nächstes ist die Effizienz der jeweiligen Tools zu betrachten. Dabei spielt
vor allem der benötigte Aufwand, das Tool ein zu setzen eine Rolle. Nebenbei
wird dabei auch noch auf die Benötigte Vorarbeit geachtet, die geleistet werden
muss um das Tool an zu wenden. Sowohl bei \lstinline{mocktest} als auch bei
\lstinline{flexmock} ist der zu betreibende Aufwand sehr gering. Beide Tools
haben ein ähnliches Interface, mit dem deskriptiv beschrieben wird, was ein
Objekt zu machen hat. Der dabei entstehende Aufwand ist sehr gering, da sich
der Code so schreibt wie ein Entwickler sich die Funktionalität vorstellt.
Beide Tools benötigen die gleiche Menge an Vorarbeit und sind demnach im Bezug
auf die Effizienz mit der sie angewendet werden können gleich auf. Lediglich
\lstinline{stubble} ist in diesem Aspekt etwas schlechter als die beiden
anderen Tools. So ist die \gls{mock}ing Funktionalität lediglich mit einem
vorgefertigten \Gls{stub} möglich, wodurch dieser erst verstanden werden muss,
bevor er effektiv eingesetzt werden kann. Zusätzlich muss jeden zu
\gls{stub}ende Objekt überschrieben werden wodurch die benötigte Vorarbeit sehr 
hoch ausfällt.

Demnach sollte im Aspekt der Effizienten Anwendung der
\gls{mock}-testing Tools \lstinline{stubble} nicht verwendet werden. Es ist dem
Entwickler überlassen ob er in diesem Aspekt das Interface von
\lstinline{mocktest} oder \lstinline{flexmock} lieber hat.
\newline

Zuletzt ist die Komplexität der \gls{mock}-testing Tools zu betrachten. Dabei
wird zum einen auf die gestellten Funktionalitäten abseits der Basis
Funktionalität geachtet und zum anderen auf den die Lesbarkeit des Entstehenden
Codes des Tools. Auch hier sind \lstinline{mocktest} und \lstinline{flexmock}
auf einem Level, was die zusätzliche Funktionalität betrifft. Jedoch hat
\lstinline{mocktest} etwas mehr zusätzliche Funktionalität geboten gegenüber
\lstinline{flexmock} (vgl. Kapitel \ref{python-tools:mocktest} und
\ref{python-tools:flexmock}). \lstinline{stubble} bietet auch in diesem Aspekt
weniger als die beiden Tools. Abseits der Basis Funktionalität wird von
\lstinline{stubble} nichts weiter geboten. Im Bezug der Lesbarkeit sieht es
ähnlich aus, \lstinline{mocktest} und \lstinline{flexmock} befinden sich auf
einem Level, während \lstinline{stubble} durch viel Code an Lesbarkeit
verliert. Im Fall \lstinline{mocktest} gegen \lstinline{flexmock} ist es wieder
dem Anwender überlassen welches der Beide Tools lesbarer ist als das andere.

Im Bezug auf Komplexität geben sich \lstinline{mocktest} und
\lstinline{flexmock} fast nichts, jedoch hat \lstinline{mocktest} durch die
zusätzliche Funktionalität hier die Nase vorne. \lstinline{stubble} bleibt auch
hier auf dem letzten Platz.
\newline

Im gesamten Vergleich ist \lstinline{mocktest} durch seine Features und sein
Interface an erster stelle, wenn es um den Einsatz von \gls{mock}-testing Tools
geht. Der zweite Platz ist schwer zu wählen, da \lstinline{flexmock} mit seiner
fehlenden Funktionalität nicht implementierte Methoden hinter
\lstinline{stubble} liegt, das jedoch auch nur bedingt diese Anforderung
erfüllt. In allen anderen Aspekten ist \lstinline{flexmock} \lstinline{stubble}
voraus. Aus diesem Grund erhält \lstinline{flexmock} den zweite Platz der
\gls{mock}-testing Tools und \lstinline{stubble} den dritten. Auf dem vierten
wäre \lstinline{doublex}, jedoch ist dieses Tool für die Anwendung von TDD
ungeeignet, weshalb es keine Platzierung erhält.
% !TeX root = ../bachelor.tex
\subsection{Kombinierung von Tools}\label{kombinierung}
Um eine verbesserte Funktionalität zu bekommen, ist es ratsam verschiedene
Tools miteinander zu kombinieren. Die in diesem Kapitel aufgezeigten
Möglichkeiten zur Kombinierung ergeben sich aus den analysierten Features der
jeweiligen Tools. Dabei wird zuerst geprüft ob die beiden, am besten geeigneten
Tools aus unit- und \gls{mock}-testing Tools miteinander verwendet werden
können, um so einen Mehrwert zu erhalten.

Bei den unit-testing Tools ist \lstinline{doctest}, das am wenigsten geeignete
Tool für die alleinige Anwendung mit TDD. Dennoch bietet das Tool Features, die
für die Anwendung von TDD nützlich sind. So kann \lstinline{doctest} als
testende Dokumentation gesehen werden. Manche Funktionen oder Methoden sind in
ihrem Funktionsumfang so gering, dass ein ganzer Test als zu viel erscheint.
Bei diesen Funktionen oder Methoden kann dann Beispielsweise ein
\lstinline{doctest} geschrieben werden, der diese kleine Funktionalität prüft.
Des weiteren kann \lstinline{doctest} als ergänzendes Tool für die
Dokumentation des Codes gesehen werden, denn die Dokumentation ist ein
wichtiger Bestandteil eines Projekts (\cite{python.org:PEP8}).
\lstinline{doctests} können dafür sorgen, dass die Dokumentation besser
verstanden wird und eindeutiger ist. Aus diesen Gründen kann
\lstinline{doctest} immer als ergänzendes Tool in jede Entwicklung aufgenommen
werden.

Da das Tool \lstinline{hypothesis} aus den \gls{fuzz} Tools, ist als alleiniger
Vertreter seiner Kategorie das empfohlene Tool, wenn \gls{fuzz} benötigt wird.
Da \lstinline{hypothesis} mit jedem anderen Tool verwendet werden kann,
empfiehlt es sich, dieses so oft wie es möglich ist ein zu setzen um eine
bessere Testabdeckung zu erreichen.

Demnach ist es ratsam sowohl \lstinline{doctest} als auch
\lstinline{hypothesis} immer ein zu setzen wo es möglich ist.
\newline

Grundsätzlich lässt sich ein jedes unit-testing Tool mit jedem 
\gls{mock}-testing Tool kombinieren. Der dadurch erhaltene Mehrwert ist meist 
beachtlich gegenüber der Verwendung der Tools alleine. Es ist jedem Entwickler 
überlassen welches Tool er mit welchem Kombiniert, dennoch hervorgehend aus der 
Analyse bietet es sich an \lstinline{pytest} zusammen mit \lstinline{mocktest} 
zu verwenden, dazu noch \lstinline{doctest} und \lstinline{hypothesis} und der 
Entwickler erhält ein Toolset, dass ihm alles an Funktionalität bietet was er 
benötigen kann. Ist es erwünscht \lstinline{unittest} statt \lstinline{pytest} 
zu verwenden, so kann dies auch getan werden, jedoch zeigt die Analyse, dass 
der Einsatz von \lstinline{pytest} dem Einsatz von \lstinline{unittest} vor zu 
ziehen ist.
    % !TeX root = ../bachelor.tex
\subsection{Zusammenfassung}\label{zusammenfassung}

Um einen schnellen gesamt Überblick zu schaffen, werden unter diesem Punkt die
einzelnen Vor- und Nachteile der Tools noch einmal zusammengefasst. Dabei wird
auf die in Kapitel \nameref{python-tools} genannten Aspekte referenziert.

Die Zusammenfassende Wertung wird mit einem Wert zwischen Eins und Zehn
bestimmt, wobei Eins der schlechteste Wert ist und Zehn der Beste. Diese Wertung
erfolgt aus den Analysen des jeweiligen Tools in den vorherigen Kapiteln. Um die
Zusammenfassung besser zu gliedern werden die Tools nicht mehr zwischen STDLIB
und externe Pakete unterschieden, sondern nach Funktionalität gegliedert (unit-,
\gls{mock}- und \gls{fuzz}-testing Tools). Zu jeder Wertung erfolgt eine kleine
Zusammenfassung der Analyse, um die Wertung nachvollziehbar zu gestalten.

%===============================================================================
% Unit-testing Tools
%===============================================================================
\subsubsection{Unit-testing Tools}\label{zusammenfassung:unit}

\paragraph{unittest}\label{zusammenfassung:unit:unittest}\mbox{}
\newline
\begin{itemize}
    \item \underline{Anwendbarkeit} (10/10):\newline
    \lstinline{unittest} bietet dank dem submodul \lstinline{unittest.mock}
    alles an Funktionalität um TDD einsetzen zu können. Da es in der STDLIB ist,
    werden keine externen Abhängigkeiten benötigt.
    
    \item \underline{Effizienz} (7/10):\newline
    Die Vorarbeit, die benötigt ist um Testen zu können hält sich in grenzen.
    Jedoch ist die Auswertung der Tests durch den Monochromen Output nicht
    sehr gut lesbar und bei vielen Fehlern unübersichtlich.
    
    \item \underline{Komplexität} (8/10):\newline
    Die von \lstinline{unittest} bereit gestellten \lstinline{assert} Methoden
    bieten dem Entwickler eine Methode für jede mögliche Behauptung. Abseits
    davon bietet \lstinline{unittest} viel Funktionalität zum aufsetzen einer
    Test Umgebung, wobei der geschriebene Code viel seiner Übersichtlichkeit
    bewahrt.
    
    \item \underline{Erweiterbarkeit} (10/10):\newline
    Viele Tools basieren oder erweitern \lstinline{unittest}, wodurch es möglich
    ist neue Funktionalität zu \lstinline{unittest} hinzu zu fügen. Die Tools,
    die \lstinline{unittest} erweitern bieten meist einen erheblichen Mehrwert
    weshalb hier 10 Punkte vergeben werden.
\end{itemize}

\paragraph{doctest}\label{zusammenfassung:unit:doctest}\mbox{}
\newline
\begin{itemize}
    \item \underline{Anwendbarkeit} (3/10):\newline
    Doctest bietet für sich alleine nicht genug Funktionalität um TDD betreiben
    zu können.
    
    \item \underline{Effizienz} (3/10):\newline
    Ein Entwickler der Python beherrscht, sollte keine Probleme haben
    \lstinline{doctest} anwenden zu können. Der Code der allerdings abseits
    der Test benötigt wird kann sehr groß ausfallen, wenn \Glspl{fixture}
    gewünscht sind. Auch der Output ist schnell unübersichtlich.
    
    \item \underline{Komplexität} (1/10):\newline
    Abgesehen der \lstinline{Option Flags} bietet \lstinline{unittest} keine
    Funktionalität um Tests zu schreiben, der Entwickler kann sich hier
    lediglich der STDLIB oder allen anderen Paketen bedienen die importiert
    werden können.
    
    \item \underline{Erweiterbarkeit} (5/10):\newline
    \lstinline{doctest} selbst ist nicht erweiterbar, kann jedoch als
    Erweiterung für andere Test Tools genutzt werden. So wird von Haus aus
    eine \lstinline{unittest} Integration bereit gestellt. Aus diesem Grund
    werden hier 5 Punkte vergeben, da das Einsatzgebiet von \lstinline{doctest}
    eher als Erweiterung gesehen werden sollte und nicht als Tool für TDD.
\end{itemize}

\paragraph{pytest}\label{zusammenfassung:unit:pytest}\mbox{}
\newline
\begin{itemize}
    \item \underline{Anwendbarkeit} (9/10):\newline
    \lstinline{pytest} bietet alles um TDD betreiben zu können, lediglich die
    der Fakt, dass es ein Externes Paket mit Sechs Abhängigkeiten ist gibt einen
    Punkt Abzug.
    
    \item \underline{Effizienz} (10/10):\newline
    Der Einstieg in \lstinline{pytest} erfolgt ohne Probleme und große
    Vorarbeit. Der Output ist gegliedert und in Farbe, wodurch er leicht zu
    lesen und zu verstehen ist.
    
    \item \underline{Komplexität} (10/10):\newline
    Durch die umfassenden Features von \lstinline{pytest} ist es dem Entwickler
    möglich viel aus seinen Tests heraus zu holen ohne neue Konstrukte lernen zu
    müssen. Der dabei entstehende Code bewahrt auch bei viel Code ein gewisses
    Maß an Lesbarkeit.
    
    \item \underline{Erweiterbarkeit} (10/10):\newline
    Die 618 (Stand 2. April 2019) verfügbaren Erweiterung rechtfertigen die Zehn
    erhaltenen Punkte.
\end{itemize}

%===============================================================================
% Mock-testing Tools
%===============================================================================
\subsubsection{Mock-testing Tools}\label{zusammenfassung:mock}

\paragraph{stubble}\label{zusammenfassung:mock:stubble}\mbox{}
\newline
\begin{itemize}
    \item \underline{Anwendbarkeit} (9/10):\newline
    Die Funktionalität, die zum \gls{stub}en benötigt wird ist mit
    \lstinline{stubble} fast komplett gegeben und die Abhängigkeiten belaufen
    sich auf nur Eine.
    
    \item \underline{Effizienz} (10/10):\newline
    Vorarbeit und viel Code werden mit \lstinline{stubble} nicht benötigt.
    
    \item \underline{Komplexität} (7/10):\newline
    \lstinline{stubble} bietet lediglich die Möglichkeit Objekte zu ersetzen,
    jedoch nicht mehr. Die Vorgefertigten \Glspl{stub} sind allerdings sehr
    nützlich. Da, keine \gls{mock}ing Funktionalität geboten wird gibt es hier
    Punktabzug.
\end{itemize}

\paragraph{mocktest}\label{zusammenfassung:mock:mocktest}\mbox{}
\newline
\begin{itemize}
    \item \underline{Anwendbarkeit} (9/10):\newline
    \lstinline{mocktest} bietet seinem Nutzer alles um Objekte zu \gls{mock}en
    oder zu \gls{stub}en. Zusätzlich erweitert \lstinline{mocktest} unittest
    wodurch es möglich ist Test Umgebungen zu erschaffen. Zusätzlich zu sich
    selbst wird Eine Abhängigkeit benötigt, für die es einen Punkt Abzug gibt.
    
    \item \underline{Effizienz} (10/10):\newline
    Da \lstinline{mocktest} deskriptiv ist, kann es ohne großen Aufwand
    verwendet werden. Der Code der dazu benötigt wird beschränkt sich auf
    ein Minimum.
    
    \item \underline{Komplexität} (10/10):\newline
    An Funktionalität fehlt es dem Tool nicht.

\end{itemize}

\paragraph{flexmock}\label{zusammenfassung:mock:flexmock}\mbox{}
\newline
\begin{itemize}
    \item \underline{Anwendbarkeit} (9/10):\newline
    Im Bezug zu TDD wird dem Anwender fast alles geboten, was benötigt wird und
    das mit nur einer Abhängigkeit, das Paket selbst.
    
    \item \underline{Effizienz} (10/10):\newline
    Die Anwendung von \lstinline{flexmock} ist deskriptiv, weshalb es für einen
    ohne größere Vorkenntnisse anwendbar ist. Der dabei zu schreibende Code ist
    Minimal.
    
    \item \underline{Komplexität} (9/10):\newline
    Das Interface von \lstinline{flexmock} ermöglicht sauberen und
    übersichtlichen Code zu schreiben und bietet dabei einiges an Funktionalität
    um zu \gls{mock}en und zu \gls{stub}en. Punktabzug gab es für die fehlende
    Möglichkeit nicht existierende Funktionen und Methoden zu erstellen.
\end{itemize}

\paragraph{doublex}\label{zusammenfassung:mock:doublex}\mbox{}
\newline
\begin{itemize}
    \item \underline{Anwendbarkeit} (X/10):\newline
    TBD.
    
    \item \underline{Effizienz} (X/10):\newline
    TBD.
    
    \item \underline{Komplexität} (X/10):\newline
    TBD.
\end{itemize}

%===============================================================================
% Fuzz-testing Tools
%===============================================================================
\subsubsection{Fuzz-testing Tools}\label{zusammenfassung:fuzz}

\paragraph{hypothesis}\label{zusammenfassung:fuzz:hypothesis}\mbox{}
\newline
\begin{itemize}
    \item \underline{Anwendbarkeit} (X/10):\newline
    TBD.
    
    \item \underline{Effizienz} (X/10):\newline
    TBD.
    
    \item \underline{Komplexität} (X/10):\newline
    TBD.
    
    \item \underline{Erweiterbarkeit} (X/10):\newline
    TBD.
\end{itemize}
    % !TeX root = ../bachelor.tex
\subsection{Vergleich der Tools}\label{vergleich}

Die in Kapitel \ref{python-tools:stdlib} und \ref{python-tools:extlib}
analysierten Tools, werden anhand der Ergebnisse aus der Anwendung des
jeweiligen Tools in diesem Kapitel verglichen. Der Vergleich wird in unit- und
\gls{mock} Tools gegliedert. Die \gls{fuzz}-testing Tools werden hier
nicht verglichen, da nur ein Tool analysiert wurde.
Am Ende dieses Kapitels wird für die jeweilige Kategorie eine Empfehlung 
ausgesprochen, die anhand der Ergebnisse des Vergleichs festgelegt wird. Anhand 
dieser Empfehlungen werden in Kapitel \ref{kombinierung} verschiedene 
Möglichkeiten aufgezeigt, die Tools untereinander zu kombinieren um mehr 
Funktionalität für eine mögliche Anwendung zu bekommen.

\subsubsection{Unit-testing Tools}\label{vergleich:unit}
Die unit-testing Tools belaufen sich auf drei Tools, \lstinline{unittest},
\lstinline{doctest} und \lstinline{pytest}. Die ersten beiden Tools sind beide
in der STDLIB während \lstinline{pytest} mit sechs weiteren Paketen betrieben
werden muss. Während \lstinline{unittest} und \lstinline{pytest} sich sehr
ähnlich sind, ist \lstinline{doctest} vorkommen anders im Bezug auf Anwendung
und gegebener Funktionalität. \lstinline{doctest} unterstützt keine Testfälle
oder Testumgebungen, stattdessen wird im \gls{docstring} der jeweiligen
Funktion oder Methode der Test geschrieben. Zwar ist es möglich diesen in eine
externe Datei aus zu lagern um so einen Zentralen Punkt zu schaffen an dem sich
alle Tests befinden (Testumgebung/Testfälle), jedoch ist hier die Datei das
Konstrukt, dass die Gliederung ermöglicht. Sowohl \lstinline{unittest} als auch
\lstinline{pytest} bieten hier die Möglichkeiten Tests in Testumgebungen
(Klassen) und Testfällen (Funktionen/Methoden) zu gliedern. \lstinline{pytest}
bietet zusätzlich auch noch die Möglichkeit diese Tests zu markieren um eine
noch feinere Gliederung zu schaffen oder um zwischen Testumgebungen, Testfälle
zu kombinieren. Im Fall der Gliederung und Verwaltung von Tests ist hier
\lstinline{pytest} sowohl \lstinline{unittest} als auch \lstinline{doctest}
voraus. Lässt man die \lstinline{unittest} Integration für \lstinline{doctest}
außen vor, so ist \lstinline{unittest} in diesem Fall der Gewinner. Wird
allerdings die Anwendung der \lstinline{unittest} Integration von
\lstinline{doctest} in Betracht gezogen, so sind \lstinline{unittest} und
\lstinline{doctest} auf einem Level, was die Gliederung und Verwaltung von
Tests betrifft.

Ein Feature, dass eine Gliederung voraussetzt sind die Fixtures. Da
\lstinline{doctest} keine Gliederung bietet sind Fixtures für dieses Tool
nicht als Feature integriert. Es ist zwar möglich Fixtures manuell zu
schreiben, jedoch ist der aufwand, der dafür betrieben werden muss zu hoch, als
dass es sich rentieren würde. Die \lstinline{unittest} Integration ermöglicht es
zwar Fixtures zu nutzen, jedoch wird diese Integration hier wieder außen vor
gelassen. \lstinline{unittest} und \lstinline{pytest} hingegen unterstützen
Testfixtures mit ihren \lstinline{setUp()} und \lstinline{tearDown()} Methoden.
Beide Tools bieten die Möglichkeit die Fixtures sowohl für jedem Testfall, jede
Testumgebung oder jedes Modul aus zu führen. \lstinline{pytest} bietet
zusätzlich noch weitere Funktionalität, sowie fertige Fixtures die dem
Entwickler zur Nutzung zur Verfügung stehen, welche in Kapitel
\ref{python-tools:pytest} beschrieben wurden und in der
\href{https://docs.pytest.org/en/latest/fixture.html}{Dokumentation}\footnote{\url{https://docs.pytest.org/en/latest/fixture.html}}
zu finden sind. Im Bezug auf Testfixtures ist also \lstinline{pytest} das Tool
mit den meisten Features, gefolgt von \lstinline{unittest}. \lstinline{doctest}
scheidet bei dieser Kategorie vollkommen aus, da die Funktionalität gänzlich
fehlt.

Als nächstes integriertes Feature der unit-testing Tools sind \Glspl{mock} und
\Glspl{stub} zu betrachten. Um für TDD in Frage zu kommen sollte ein
unit-testing Tool \Glspl{stub} und \Glspl{mock} unterstützen, eine ausführliche
Erklärung dazu ist in Kapitel \ref{python-tools} zu finden. Wie in Kapitel
\ref{python-tools:doctest} aus der Analyse von \lstinline{doctest} hervorgeht,
gibt es weder die Möglichkeit einen \Gls{mock} noch einen \Gls{stub} auf zu
setzen, weshalb \lstinline{doctest} in dieser Kategorie als unbrauchbar gesehen
werden kann. \lstinline{unittest} als Modul selbst, besitzt keine \gls{mock}ing
Funktionalität, jedoch ist es möglich, dass Submodul \lstinline{unittest.mock}
dafür zu verwenden. \lstinline{unittest.mock} ist zwar eine Erweiterung, jedoch
in Form eines Submoduls, weshalb es zu \lstinline{unittest} dazu gehört und in
diesem Vergleich beachtet wird. Damit unterstützt \lstinline{unittest} sowohl
das \gls{mock}en als auch das \gls{stub}en. \lstinline{pytest} hingegen verfügt
in seiner Basisausführung lediglich über die Möglichkeit des \gls{stub}ens,
nicht jedoch über die Möglichkeit ein Objekt zu \gls{mock}en. In diesem Falls
ist also \lstinline{unittest} der klare Gewinner, gefolgt von
\lstinline{pytest}.
\newline

Im Bezug auf die Effizienz der unit-testing Tools lässt sich folgendes
feststellen. \lstinline{doctest} ist zwar für jeden Entwickler einfach und
verständlich anwendbar, jedoch ist die Vorarbeit um Tests auf zu setzen im
Vergleich zu \lstinline{unittest} und \lstinline{pytest} höher. Sowohl
\lstinline{unittest} als auch \lstinline{pytest} benötigen das gleiche Maß an
Vorarbeit um Tests auf zu setzen, wobei der Aufwand das jeweilige Tool ein zu
setzen bei \lstinline{unittest} etwas höher ist, da es so viele verschiedene
\lstinline{assert} Methoden zu merken gibt. Die Auswertung der Tests erfolgt
bei allen drei unit-testing Tools auf dem Terminal. \lstinline{doctest} und
\lstinline{unittest} haben beide einen Output der voneinander fast nicht zu
trennen ist. Beide Tools haben den Nachteil gegenüber \lstinline{pytest}, dass
ihr Output nicht farbig sonder Monochrom ist.  Aus diesem Grund ist 
\lstinline{pytest}, \lstinline{unittest} und \lstinline{doctest} im Bezug auf
Effiziente Anwendung vor zu ziehen, wobei \lstinline{unittest} noch etwas 
besser als \lstinline{doctest} ist.
\newline

Der Aspekt der Komplexität ist in zwei unter Kategorien gegliedert, einerseits
ist da die Fülle an Funktionen die geboten werden und Komplexität liefern und
zum anderen ist da die Menge an Code die zu schreiben ist und die Anwendung und
Wartung Komplexer gestaltet. Im Falle der ersten Kategorie, erweist sich
\lstinline{doctest} als nicht all zu komplex. Fast die Gesamte gebotene
Funktionalität bezieht sich auf die Ausführung, jedoch nicht das schreiben der
Tests, welches alleine durch die Fähigkeit Python zu schreiben bestimmt wird.
\lstinline{unittest} und \lstinline{pytest} hingegen bieten dem Entwickler
einiges an Funktionalität um Tests zu schreiben und aus zu führen. Auch hier
kann \lstinline{unittest} nur mithalten, wenn das Submodul
\lstinline{unittest.mock} mit in die Wertung aufgenommen wird, da dieses
weitere Funktionalität bietet. \lstinline{unittest} und \lstinline{pytest}
haben viele gemeinsame Features, wovon lediglich das \gls{mock}ing in
\lstinline{pytest} nicht ohne Erweiterungen verfügbar ist. Ansonsten bietet
\lstinline{pytest} viel zusätzliche Funktionalität um Tests feiner ein zu
stellen oder um irgendwelche edge cases ab zu decken. In dieser Kategorie ist
\lstinline{pytest} also etwas besser als \lstinline{unittest} und beide Tools
sind \lstinline{doctest} bei weitem überlegen im Bezug auf die gebotenen
Funktionalitäten.

In der zweiten Kategorie, der Komplexität im Bezug auf Code Lesbarkeit und
Umfang erweisen sich \lstinline{unittest} und \lstinline{pytest} als gleichauf.
Beide Tools haben einen ähnlichen Aufbau, mit den der Code geschrieben wird
(vgl. Listing \ref{listing:unittest:advanced} und
\ref{listing:pytest:advanced}). Die Lesbarkeit ist bei beiden Tools gleich gut
und die Menge an Code die benötigt wird gleicht sich auch. Hier ist es dem
Entwickler überlassen zu entscheiden, welches Tool den schöneren Code schreibt.
\lstinline{doctest} hingegen, produziert viel Code, der mit zunehmender
Funktionalität an Lesbarkeit verliert.Im Bezug auf die Komplexität der Tools
ist \lstinline{pytest} also knapp vor \lstinline{unittest}, gefolgt von
\lstinline{doctest}.
\newline

Als letztes ist die Erweiterbarkeit der unit-testing Tools zu vergleichen.
Beginnend mit \lstinline{doctest}, lässt sich sagen, das es zwei sehr gute
Erweiterungen gibt, die das Tool erheblich verbessern. Diese Erweiterungen sind
die Integration in \lstinline{unittest} und \lstinline{pytest}, erstere wird
von der Python Sprache selbst geboten und ist in \lstinline{unittest} verfügbar
und zweite wird von \lstinline{pytest} geboten und ist auch im Basis Paket
enthalten. Abgesehen von diesen zwei gibt es noch weitere Integrationen in
andere Toolsets, jedoch wurden diese hier aus verschiedenen Gründen nicht
behandelt und werden deshalb nicht erwähnt. Das Modul \lstinline{unittest}
verfügt über einige Erweiterungen welche neue Funktionalität hinzu fügen oder
alte Funktionalität verbessern. Das Tool mit den meisten Erweiterungen ist
\lstinline{pytest}, welches derzeit 618 (Stand: 2. April
2019\footnote{\url{http://plugincompat.herokuapp.com/}}) Erweiterungen
bietet, wovon die meisten sehr nützliche Funktionalität bieten.
\newline

Im Vergleich der unit-testing Tools ergibt sich also ein klarer Gewinner,
\lstinline{pytest}. Das Tool ist in fast allen Analysierten Aspekten die Beste
Wahl für den Einsatz mit TDD. Jedoch bietet \lstinline{unittest} für
Entwickler, die keine Externen Abhängigkeiten möchten eine alternative, die
durchaus mit \lstinline{pytest} mithalten kann.

\subsubsection{Mock-testing Tools}\label{vergleich:mock}
Die analysierten \gls{mock}-testing Tools, sind \lstinline{stubble},
\lstinline{mocktest}, \lstinline{flexmock} und \lstinline{doublex}. Alle vier
bieten verschiedene Möglichkeiten \Glspl{mock} und/oder \Glspl{stub} ein zu
setzen. Da es sich hier um Tools handelt, deren Funktionalität sich um das
\gls{mock}en von Objekten dreht, wird der erste Analysierte Aspekt, die
Anwendbarkeit, auf die gebotene Auswahl an Funktionalität im Bezug auf
\gls{mock}en und \gls{stub}en betrachtet.

Das einzige Tool, das bei der Anwendung auf den zu Testenden Code keine
Probleme mit sich barg, ist \lstinline{mocktest}. Das Tool bietet im Aspekt der
Anwendbarkeit sämtliche Features um effektiv \Glspl{mock} und \Glspl{stub} ein
zu setzen. \lstinline{mocktest} bietet seinem Anwender eine Deskriptive schreib
weise, mit der es problemlos möglich ist Tests zu schreiben. Des weiteren wird
eine Testisolation geboten, mit deren Hilfe es möglich ist Test besser
voneinander ab zu schotten.Zusätzlich benötigt das Tool keine weiteren
Abhängigkeiten zu sich selbst.
Direkt danach wäre \lstinline{stubble} das  nächst beste Tool im Aspekt der
Anwendbarkeit. So bietet das Tool zwar alles um  es für TDD zu verwenden,
jedoch ist das \gls{stub}en einer nicht existierenden Methode nur unter
Einschränkungen Möglich. Die Hauptfeatures von \lstinline{stubble} sind dabei
die vorgefertigten \Glspl{stub}, welche nützliche Features für den Entwickler
bieten. Um einen \Gls{stub} zu erstellen muss die zu \gls{stub}ende Klasse mit
\lstinline{stubble} neu geschrieben werden.Zusätzlich benötigt
\lstinline{stubble} eine weitere Abhängigkeit zu sich selbst.
An dritter Stelle ist das Tool \lstinline{flexmock} welches alles bietet,
abgesehen vom ersetzen einer nicht geschriebenen Methode, um TDD an zu wenden
(vgl. Listing \ref{listing:flexmock:exception}). Wäre dieses Manko nicht, würde
das Tool im Aspekt der Anwendbarkeit sehr weit oben stehen. Flexmock vereint
die \Glspl{mock} und \Glspl{stub} in einer Klasse wodurch die gesamte
Funktionalität mit einem Objekt geboten wird. Die schreib weise der Tests ist
hierbei, wie bei \lstinline{mocktest}, deskriptiv. Auch die benötigten
Abhängigkeiten von \lstinline{flexmock} sprechen für das Tool, da keine
weiteren zu sich selbst benötigt werden. Als letztes Tool im Aspekt der
Anwendbarkeit ist \lstinline{doublex} zu betrachten. Das Tool bietet zwar
einiges an Funktionalität um zu \gls{mock}en oder zu \gls{stub}en, jedoch ist
die Anwendung dieser Features nicht nutzbar für TDD. Es werden vier
verschiedene Interfaces geboten, deren Funktionalität sich überschneidet und
die jeweils ein Objekt hervorbringen, dass ein Duplikat vom Ziel Objekt ist.
Die daraus resultierenden Anwendungszwecke bieten wenig Nutzbarkeit für TDD.
Hinzu kommt, dass \lstinline{doublex} drei weitere Abhängigkeiten benötigt und
anscheinend nicht mehr aktiv weiter entwickelt wird.

Demnach ist im Aspekt der Anwendbarkeit \lstinline{mocktest} also
\lstinline{stubble} und \lstinline{flexmock} vor zu ziehen. \lstinline{doublex}
fällt aufgrund dessen, dass es nicht für TDD geeignet ist und derzeit nicht
aktiv weiter entwickelt wird komplett raus. Aus diesem Grund wird dieses Tool
im weiteren verlauf des Vergleichs nicht mehr berücksichtigt.
\newline

Als nächstes ist die Effizienz der jeweiligen Tools zu betrachten. Dabei spielt
vor allem der benötigte Aufwand, das Tool ein zu setzen eine Rolle. Nebenbei
wird dabei auch noch auf die Benötigte Vorarbeit geachtet, die geleistet werden
muss um das Tool an zu wenden. Sowohl bei \lstinline{mocktest} als auch bei
\lstinline{flexmock} ist der zu betreibende Aufwand sehr gering. Beide Tools
haben ein ähnliches Interface, mit dem deskriptiv beschrieben wird, was ein
Objekt zu machen hat. Der dabei entstehende Aufwand ist sehr gering, da sich
der Code so schreibt wie ein Entwickler sich die Funktionalität vorstellt.
Beide Tools benötigen die gleiche Menge an Vorarbeit und sind demnach im Bezug
auf die Effizienz mit der sie angewendet werden können gleich auf. Lediglich
\lstinline{stubble} ist in diesem Aspekt etwas schlechter als die beiden
anderen Tools. So ist die \gls{mock}ing Funktionalität lediglich mit einem
vorgefertigten \Gls{stub} möglich, wodurch dieser erst verstanden werden muss,
bevor er effektiv eingesetzt werden kann. Zusätzlich muss jeden zu
\gls{stub}ende Objekt überschrieben werden wodurch die benötigte Vorarbeit sehr 
hoch ausfällt.

Demnach sollte im Aspekt der Effizienten Anwendung der
\gls{mock}-testing Tools \lstinline{stubble} nicht verwendet werden. Es ist dem
Entwickler überlassen ob er in diesem Aspekt das Interface von
\lstinline{mocktest} oder \lstinline{flexmock} lieber hat.
\newline

Zuletzt ist die Komplexität der \gls{mock}-testing Tools zu betrachten. Dabei
wird zum einen auf die gestellten Funktionalitäten abseits der Basis
Funktionalität geachtet und zum anderen auf den die Lesbarkeit des Entstehenden
Codes des Tools. Auch hier sind \lstinline{mocktest} und \lstinline{flexmock}
auf einem Level, was die zusätzliche Funktionalität betrifft. Jedoch hat
\lstinline{mocktest} etwas mehr zusätzliche Funktionalität geboten gegenüber
\lstinline{flexmock} (vgl. Kapitel \ref{python-tools:mocktest} und
\ref{python-tools:flexmock}). \lstinline{stubble} bietet auch in diesem Aspekt
weniger als die beiden Tools. Abseits der Basis Funktionalität wird von
\lstinline{stubble} nichts weiter geboten. Im Bezug der Lesbarkeit sieht es
ähnlich aus, \lstinline{mocktest} und \lstinline{flexmock} befinden sich auf
einem Level, während \lstinline{stubble} durch viel Code an Lesbarkeit
verliert. Im Fall \lstinline{mocktest} gegen \lstinline{flexmock} ist es wieder
dem Anwender überlassen welches der Beide Tools lesbarer ist als das andere.

Im Bezug auf Komplexität geben sich \lstinline{mocktest} und
\lstinline{flexmock} fast nichts, jedoch hat \lstinline{mocktest} durch die
zusätzliche Funktionalität hier die Nase vorne. \lstinline{stubble} bleibt auch
hier auf dem letzten Platz.
\newline

Im gesamten Vergleich ist \lstinline{mocktest} durch seine Features und sein
Interface an erster stelle, wenn es um den Einsatz von \gls{mock}-testing Tools
geht. Der zweite Platz ist schwer zu wählen, da \lstinline{flexmock} mit seiner
fehlenden Funktionalität nicht implementierte Methoden hinter
\lstinline{stubble} liegt, das jedoch auch nur bedingt diese Anforderung
erfüllt. In allen anderen Aspekten ist \lstinline{flexmock} \lstinline{stubble}
voraus. Aus diesem Grund erhält \lstinline{flexmock} den zweite Platz der
\gls{mock}-testing Tools und \lstinline{stubble} den dritten. Auf dem vierten
wäre \lstinline{doublex}, jedoch ist dieses Tool für die Anwendung von TDD
ungeeignet, weshalb es keine Platzierung erhält.
    % !TeX root = ../bachelor.tex
\subsection{Kombinierung von Tools}\label{kombinierung}
Um eine verbesserte Funktionalität zu bekommen, ist es ratsam verschiedene
Tools miteinander zu kombinieren. Die in diesem Kapitel aufgezeigten
Möglichkeiten zur Kombinierung ergeben sich aus den analysierten Features der
jeweiligen Tools. Dabei wird zuerst geprüft ob die beiden, am besten geeigneten
Tools aus unit- und \gls{mock}-testing Tools miteinander verwendet werden
können, um so einen Mehrwert zu erhalten.

Bei den unit-testing Tools ist \lstinline{doctest}, das am wenigsten geeignete
Tool für die alleinige Anwendung mit TDD. Dennoch bietet das Tool Features, die
für die Anwendung von TDD nützlich sind. So kann \lstinline{doctest} als
testende Dokumentation gesehen werden. Manche Funktionen oder Methoden sind in
ihrem Funktionsumfang so gering, dass ein ganzer Test als zu viel erscheint.
Bei diesen Funktionen oder Methoden kann dann Beispielsweise ein
\lstinline{doctest} geschrieben werden, der diese kleine Funktionalität prüft.
Des weiteren kann \lstinline{doctest} als ergänzendes Tool für die
Dokumentation des Codes gesehen werden, denn die Dokumentation ist ein
wichtiger Bestandteil eines Projekts (\cite{python.org:PEP8}).
\lstinline{doctests} können dafür sorgen, dass die Dokumentation besser
verstanden wird und eindeutiger ist. Aus diesen Gründen kann
\lstinline{doctest} immer als ergänzendes Tool in jede Entwicklung aufgenommen
werden.

Da das Tool \lstinline{hypothesis} aus den \gls{fuzz} Tools, ist als alleiniger
Vertreter seiner Kategorie das empfohlene Tool, wenn \gls{fuzz} benötigt wird.
Da \lstinline{hypothesis} mit jedem anderen Tool verwendet werden kann,
empfiehlt es sich, dieses so oft wie es möglich ist ein zu setzen um eine
bessere Testabdeckung zu erreichen.

Demnach ist es ratsam sowohl \lstinline{doctest} als auch
\lstinline{hypothesis} immer ein zu setzen wo es möglich ist.
\newline

Grundsätzlich lässt sich ein jedes unit-testing Tool mit jedem 
\gls{mock}-testing Tool kombinieren. Der dadurch erhaltene Mehrwert ist meist 
beachtlich gegenüber der Verwendung der Tools alleine. Es ist jedem Entwickler 
überlassen welches Tool er mit welchem Kombiniert, dennoch hervorgehend aus der 
Analyse bietet es sich an \lstinline{pytest} zusammen mit \lstinline{mocktest} 
zu verwenden, dazu noch \lstinline{doctest} und \lstinline{hypothesis} und der 
Entwickler erhält ein Toolset, dass ihm alles an Funktionalität bietet was er 
benötigen kann. Ist es erwünscht \lstinline{unittest} statt \lstinline{pytest} 
zu verwenden, so kann dies auch getan werden, jedoch zeigt die Analyse, dass 
der Einsatz von \lstinline{pytest} dem Einsatz von \lstinline{unittest} vor zu 
ziehen ist.
    % !TeX root = ../bachelor.tex
\section{Diskussion}\label{diskussion}
Um für die Anwendung von TDD die passenden Tools zu finden, wurden verschiedene
Tool der Programmiersprache Python analysiert und auf ihre Tauglichkeit im
Bezug auf TDD geprüft. Daraus entstand eine Empfehlung für vier Tools die
gemeinsam für TDD genutzt werden können, um so die optimale Funktionalität zum
schreiben von Tests, sowie zu einer optimalen Testabdeckung zu gewähren. Die 
dabei erhaltenen Ergebnisse zeigen auf, dass ein Tool alleine zwar ausreichend 
sein kann für die Anwendung von TDD, jedoch ist die Verwendung von mehreren 
Tools zusammen für den Entwickler und für das Produkt besser.

Im Rahmen dieser Arbeit wurden lediglich Tools analysiert, deren letzte
Aktivität nicht weiter als ein Jahr zurück liegt. Tools die eine Erweiterung zu
einem anderen Tool darstellen wurden nicht behandelt, da dies den Rahmen der
Arbeit gesprengt hätte. Aus diesem Grund kann es sein, dass ein Beliebtes oder
Bekanntes Tool in dieser Arbeit nicht behandelt wurde. Das bedeutet nicht, dass
diese Tool zwangsläufig nicht geeignet ist für TDD. Die Ergebnisse dieser
Arbeit sind also nicht repräsentativ für alle Tools die es für die
Programmiersprache Python gibt und würden eventuell anders ausfallen, wenn
andere Tools behandelt worden würden.

Im Bezug auf die Forschungsfrage, welches aktuelle Python Test-Tool das beste 
ist, lässt sich folgendes sagen: Viele Tools sind in ihrem Funktionsumfang sehr 
ähnlich, allerdings stellten sich \lstinline{pytest} bei den unit-testing-, 
\lstinline{mocktest} bei den mock-testing- und \lstinline{hypothesis} bei den 
fuzz-testing Tools als die beste Wahl heraus. Dennoch ergab sich, dass 
\lstinline{doctest} immer als ergänzendes Tool verwendet werden kann, egal 
welche anderen Tools verwendet werden.

Diese Arbeit hat sich jedoch lediglich mit einer begrenzten Anzahl von Tools 
beschäftigt und hat diese Anhand eines Spezifischen Beispiels getestet. Jeder 
Entwickler oder Software-Architekt sollte sich vorher noch einmal genaustens 
darüber informieren welches Tool für den geplanten Anwendungszweck das beste 
ist, da manche Tools abhängig ihrer Anwendung vielleicht eine bessere Wahl 
darstellen.
    % !TeX root = ../bachelor.tex
\section{Fazit}\label{fazit}

FAZIT

    % !TeX root = ../bachelor.tex
\section{Nachwort}\label{nachwort}

Nachwort


    \onecolumn

    % einfacher Zeilenabstand
    \singlespacing

    % Literaturliste soll im Inhaltsverzeichnis auftauchen
    \newpage
    \addcontentsline{toc}{section}{Literaturverzeichnis}

    % Literaturverzeichnis anzeigen
    \renewcommand\refname{Literaturverzeichnis}
    \printbibliography

    % 1,5 facher Zeilenabstand
    \onehalfspacing

    \newpage

    % Eidesstattliche Erklärung
    \addcontentsline{toc}{section}{Eidesstattliche Erklärung}
    % !TeX root = bachelor.tex
\newpage

\thispagestyle{empty}
\addcontentsline{toc}{section}{Eidesstattliche Erklärung}

\begin{verbatim}

\end{verbatim}

\begin{LARGE}
\begin{center}
	\textbf{Eidesstattliche Erklärung zur Abschlussarbeit}
\end{center}
\end{LARGE}
\begin{verbatim}


\end{verbatim}
Hiermit versichere ich, die eingereichte Abschlussarbeit selbständig verfasst und
keine andere als die von mir angegebenen Quellen und Hilfsmittel benutzt zu haben.
Wörtlich oder inhaltlich verwendete Quellen wurden entsprechend den anerkannten
Regeln wissenschaftlichen Arbeitens zitiert. Ich erkläre weiterhin, dass die
vorliegende Arbeit noch nicht anderweitig als Abschlussarbeit eingereicht wurde.
\\
\\
Das Merkblatt zum Täuschungsverbot im Prüfungsverfahren der Hochschule
Augsburg habe ich gelesen und zur Kenntnis genommen. Ich versichere, dass die
von mir abgegebene Arbeit keinerlei Plagiate, Texte oder Bilder umfasst, die durch
von mir beauftragte Dritte erstellt wurden.
\begin{verbatim}



\end{verbatim}


\begin{displaymath}
% use packages: array
\begin{array}{ll}
Unterschrift:~~~~~~~~~~~~~~~~~~~~~~~~~~~~~~~~~~~~~~~~~~
& Ort, Datum:~~~~~~~~~~~~~~~~~~~~~~~~~~~~~~~~~~~~~~~~~~
\end{array}
\end{displaymath}

    % leere Abschlussseite
    \newpage
    \thispagestyle{empty} % erzeugt Seite ohne Kopf- / Fusszeile
\end{document}
